% ==================================================================
% id:       func-014
% topic:    פונקציות
% subtopic: תחום הגדרה ותכונות של פונקציות
% solution: full
% ==================================================================
\textbf{בדקו האם הפונקציות הבאות מחזוריות
, ואם כן, מצאו מחזור:}
\begin{enumerate}
\item $f(x) =\cos(x)+\sin(x) $
\begin{solbox}
הפונקציה מחזורית. למשל $T=2\pi$ הוא מחזור, כי
\[
.f(x+2\pi)=\cos(x+2\pi)+\sin(x+2\pi)=\cos x+\sin x=f(x)
\]
\end{solbox}

\item $f(x) = \cos^2(x) $
\begin{solbox}
הפונקציה מחזורית, ו-$T=\pi$ הוא מחזור:
\[
.\cos^2(x+\pi)=(-\cos x)^2=\cos^2x
\]
כמובן שגם
$2\pi$
יהיה מחזור שלה.
\end{solbox}

\item $f(x) = 3.5 $
\begin{solbox}
הפונקציה קבועה, ולכן כל $T>0$ הוא מחזור שלה.
\end{solbox}

\item $f(x) = \sin(2x) $
\begin{solbox}
הפונקציה מחזורית, ו-$T=\pi$ הוא מחזור:
\[
.\sin(2(x+\pi))=\sin(2x+2\pi)=\sin(2x)
\]
\end{solbox}

\item פונקצית דיריכלה: $D(x) = \case{cc}{
       1 & \text{\R{ רציונאלי }} x\\ 0 & \text{\R{ אי-רציונאלי }} x}$
\hint{בדקו עבור אילו מספרים $T$ מובטח ש-$x$ ו-$x+T$ הם מאותו סוג (רציונלי או אי-רציונלי) לכל $x$.}
\begin{solbox}
כל מספר רציונלי
$q\in\mathbb{Q}$
הוא מחזור של פונקציית דיריכלה: אם
$x$
רציונלי אז גם
$x+q$
רציונלי,ולכן
$$.D(x)=1=D(x+q)$$
ואם
$x$
אי-רציונלי אז גם
$x+q$
אי-רציונלי ולכן
$$.D(x)=0=D(x+q)$$
לכן
$D(x+q)=D(x)$
לכל
$x$.
מספר
$r$
אי-רציונלי אינו מחזור, כי למשל עבור
$x=0$
הערכים
$D(0)$ ו-$D(r)$ שונים.
\end{solbox}
\item $f(x) = \cosh(x) $
\hint{אם נראה לכם שהפונקציה מחזורית, מצאו מחזור; אחרת חפשו תכונה של $\cosh$ שסותרת את $f(x+T)=f(x)$ לכל $x$.}
\begin{solbox}
קל לראות שהפונקציה עולה ממש לכל איקס חיובי ולכן לא יכולה להיות מחזורית
\end{solbox}
\item $f(x) = a^x $
הדרכה: הפרידו למקרים
\hint{הפרידו בין המקרה $a=1$ לשאר הערכים של $a$, ובכל מקרה בדקו את המונוטוניות של הפונקציה.}
\begin{solbox}
עבור
$a=1$
ברור כי קבועה ולכן מחזורית
\\
במקרה והבסיס $a$
שונה מ-1 הפונקציה מונוטונית ממש ולכן לא מחזורית.
\end{solbox}
\end{enumerate}
